\chapter{2023年全国乙卷（文）}

\section{单选题}

\begin{problem}
\(\left|2+\i^2+2\i^3\right|=\)（\quad）
\choices{\(1\)}{\(2\)}{\(\sqrt5\)}{\(5\)}
\end{problem}

\begin{answer}C\end{answer}

\begin{solution}\(2+\i^2+2\i^3=1-2\i\)，所以 \(\left|1-2\i\right|=\sqrt5\). 故选 C.
\end{solution}

\begin{problem}
设全集 \(U=\{0,1,2,4,6,8\}\)，集合 \(M=\{0,4,6\}\)，\(N=\{0,1,6\}\)，则 \(M\cup(U\setminus N)\) 等于（\quad）
\choices{\(\{0,2,4,6,8\}\)}{\(\{0,1,4,6,8\}\)}{\(\{1,2,4,6,8\}\)}{\(U\)}
\end{problem}

\begin{answer}A\end{answer}

\begin{solution}由 \(U\setminus N=\{2,4,8\}\)，得 \(M\cup(U\setminus N)=\{0,2,4,6,8\}\). 故选 A.
\end{solution}

\begin{problem}
如图，网格纸上绘制的一个零件的三视图中，网格小正方形边长为 \(1\)，则该零件的表面积为（\quad）
\begin{center}\bitmapfigure[max width=6.0cm]{img/2023/national_paper_B_liberal/q3_fig1.png}\end{center}
\choices{\(24\)}{\(26\)}{\(28\)}{\(30\)}
\end{problem}

\begin{answer}D\end{answer}

\begin{solution}该零件可看作长、宽、高分别为 \(2\)，\(2\)，\(3\) 的长方体去掉两个 \(1\times1\) 小正方体，故表面积为 \(2(2\cdot2)+4(2\cdot3)-2=30\). 故选 D.
\end{solution}

\begin{problem}
记 \(\triangle ABC\) 的内角 \(A\)，\(B\)，\(C\) 的对边分别为 \(a\)，\(b\)，\(c\)，若 \(a\cos B-b\cos A=c\)，且 \(C=\frac\pi5\)，则 \(B=\)（\quad）
\choices{\(\frac\pi{10}\)}{\(\frac\pi5\)}{\(\frac{3\pi}{10}\)}{\(\frac{2\pi}5\)}
\end{problem}

\begin{answer}C\end{answer}

\begin{solution}由正弦定理，题设等式化为 \(\sin(A-B)=\sin C\). 因 \(A+B=\pi-C\)，可得 \(A=\frac\pi2\)，从而 \(B=\frac\pi2-C=\frac{3\pi}{10}\). 故选 C.
\end{solution}

\begin{problem}
已知 \(f(x)=\dfrac{xe^x}{e^{ax}-1}\) 是偶函数，则 \(a=\)（\quad）
\choices{\(-2\)}{\(-1\)}{\(1\)}{\(2\)}
\end{problem}

\begin{answer}D\end{answer}

\begin{solution}偶函数要求 \(f(x)=f(-x)\). 对 \(x\ne0\) 化简得 \(e^x=e^{(a-1)x}\)，故 \(a=2\). 故选 D.
\end{solution}

\begin{problem}
已知正方形 \(ABCD\) 的边长为 \(2\)，\(E\) 为 \(AB\) 的中点，则 \(\overrightarrow{EC}\cdot\overrightarrow{ED}=\)（\quad）
\choices{\(\sqrt5\)}{\(3\)}{\(2\sqrt5\)}{\(5\)}
\end{problem}

\begin{answer}B\end{answer}

\begin{solution}取 \(A(0,0)\)，\(B(2,0)\)，\(C(2,2)\)，\(D(0,2)\)，\(E(1,0)\)，则 \(\overrightarrow{EC}=(1,2)\)，\(\overrightarrow{ED}=(-1,2)\)，所以 \(\overrightarrow{EC}\cdot\overrightarrow{ED}=3\). 故选 B.
\end{solution}

\begin{problem}
设 \(O\) 为平面直角坐标系原点，在区域 \(\{(x,y)\mid1\le x^2+y^2\le4\}\) 内随机取一点 A，则直线 \(OA\) 的倾斜角不大于 \(\frac\pi4\) 的概率为（\quad）
\choices{\(\frac18\)}{\(\frac16\)}{\(\frac14\)}{\(\frac12\)}
\end{problem}

\begin{answer}C\end{answer}

\begin{solution}区域为圆环，圆心角与面积成正比. 直线的倾斜角在 \([0,\pi)\) 内均匀，故所求概率为 \(\dfrac{\pi/4}{\pi}=\frac14\). 故选 C.
\end{solution}

\begin{problem}
若函数 \(f(x)=x^3+ax+2\) 有 \(3\) 个零点，则 \(a\) 的取值范围是（\quad）
\choices{\((-\infty,-2)\)}{\((-\infty,-3)\)}{\((-4,-1)\)}{\((-3,0)\)}
\end{problem}

\begin{answer}B\end{answer}

\begin{solution}要有三个零点，需 \(a<0\) 且极大值大于零、极小值小于零. 令 \(f'(x)=3x^2+a\)，代入两个极值点可得 \(a<-3\). 故选 B.
\end{solution}

\begin{problem}
某学校举办作文比赛，共设 \(6\) 个主题，每位参赛同学从中随机抽取一个主题，则甲、乙两位参赛同学抽到不同主题的概率为（\quad）
\choices{\(\frac56\)}{\(\frac23\)}{\(\frac12\)}{\(\frac13\)}
\end{problem}

\begin{answer}A\end{answer}

\begin{solution}总情况数为 \(6^2\)，抽到不同主题的情况数为 \(6\cdot5\)，故概率为 \(\frac{30}{36}=\frac56\). 故选 A.
\end{solution}

\begin{problem}
已知函数 \(f(x)=\sin(\omega x+\varphi)\) 在区间 \(\left(\frac\pi6,\frac{2\pi}3\right)\) 上单调递增，直线 \(x=\frac\pi6\) 和 \(x=\frac{2\pi}3\) 为函数 \(y=f(x)\) 的图象的两条对称轴，则 \(f\left(-\frac{5\pi}{12}\right)=\)（\quad）
\choices{\(-\frac{\sqrt3}{2}\)}{\(-\frac12\)}{\(\frac12\)}{\(\frac{\sqrt3}{2}\)}
\end{problem}

\begin{answer}D\end{answer}

\begin{solution}两条相邻对称轴间距为 \(\frac\pi\omega\)，故 \(\omega=2\). 由单调性，\(f(\frac\pi6)=-1\)，于是 \(f(-\frac{5\pi}{12})=\sin\left(\frac{3\pi}2+2\left(-\frac{5\pi}{12}-\frac\pi6\right)\right)=\frac{\sqrt3}{2}\). 故选 D.
\end{solution}

\begin{problem}
已知 \(x\)，\(y\) 满足 \(x^2+y^2-4x-2y-4=0\)，则 \(x-y\) 的最大值为（\quad）
\choices{\(1+\frac{3\sqrt2}{2}\)}{\(4\)}{\(1+3\sqrt2\)}{\(7\)}
\end{problem}

\begin{answer}C\end{answer}

\begin{solution}圆心为 \((2,1)\)，半径为 \(3\). 因此 \(x-y\) 的最大值为 \(2-1+3\sqrt2=1+3\sqrt2\). 故选 C.
\end{solution}

\begin{problem}
设 \(A\)，\(B\) 是双曲线 \(x^2-\dfrac{y^2}{9}=1\) 上两点，下列四个点中可以为线段 \(AB\) 中点的是（\quad）
\choices{\((1,1)\)}{\((-1,2)\)}{\((1,3)\)}{\((-1,-4)\)}
\end{problem}

\begin{answer}D\end{answer}

\begin{solution}设中点为 \(M(u,v)\)，弦方向向量可取 \((v,9u)\). 由双曲线方程，两端点存在的必要条件为 \((v^2-9u^2)\bigl(9-9u^2+v^2\bigr)>0\). 逐项检验可知只有 \((-1,-4)\) 满足条件. 故选 D.
\end{solution}

\section{填空题}

\begin{problem}
已知点 \(A(1,\sqrt5)\) 在抛物线 \(C:y^2=2px\) 上，则 \(A\) 到 \(C\) 的准线的距离为\fillinblank{}.
\end{problem}

\begin{answer}\(\dfrac94\)\end{answer}

\begin{solution}由 \(5=2p\) 得 \(p=\frac52\)，准线为 \(x=-\frac54\)，故点 \(A\) 到准线的距离为 \(1+\frac54=\frac94\).
\end{solution}

\begin{problem}
若 \(\theta\in\left(0,\frac\pi2\right)\)，且 \(\tan\theta=\frac13\)，则 \(\sin\theta-\cos\theta=\)\fillinblank{}.
\end{problem}

\begin{answer}\(-\dfrac{\sqrt{10}}5\)\end{answer}

\begin{solution}由 \(\tan\theta=\frac13\) 及锐角条件，得 \(\sin\theta=\frac1{\sqrt{10}}\)，\(\cos\theta=\frac3{\sqrt{10}}\)，所以 \(\sin\theta-\cos\theta=-\frac{\sqrt{10}}5\).
\end{solution}

\begin{problem}
若 \(x\)，\(y\) 满足约束条件 \(\begin{cases}
x-3y\le-1,\\
x+2y\le9,\\
3x+y\ge7,
\end{cases}\) 则 \(z=2x-y\) 的最大值为\fillinblank{}.
\end{problem}

\begin{answer}\(8\)\end{answer}

\begin{solution}可行域顶点包括 \((2,1)\)，\((1,4)\)，\((5,2)\)，在这些点上计算 \(z\) 得 \(3\)，\(-2\)，\(8\)，故最大值为 \(8\).
\end{solution}

\begin{problem}
已知点 \(S\)，\(A\)，\(B\)，\(C\) 均在半径为 \(2\) 的球面上，\(\triangle ABC\) 是边长为 \(3\) 的等边三角形，且 \(SA\perp\) 平面 \(ABC\)，则 \(SA=\)\fillinblank{}.
\end{problem}

\begin{answer}\(2\)\end{answer}

\begin{solution}设 \(O\) 为正三角形 \(ABC\) 的中心，则 \(OA=\frac3{\sqrt3}=\sqrt3\). 球心在过 \(O\) 的平面 \(ABC\) 的垂线上，球心到平面的距离为 \(\sqrt{2^2-OA^2}=1\). 由 \(SA\perp\) 平面 \(ABC\) 且 \(S\) 在球面上，得 \(SA=2\).
\end{solution}

\section{解答题}

\begin{problem}
某厂为比较甲、乙两种工艺对橡胶产品伸缩率的处理效应，进行 \(10\) 次配对试验，每次配对试验选用材质相同的两个橡胶产品，随机地选其中一个用甲工艺处理，另一个用乙工艺处理，测量处理后的橡胶产品的伸缩率. 甲、乙两种工艺处理后的橡胶产品的伸缩率分别记为 \(x_i\)，\(y_i\)（\(i=1\)，\(2\)，\(\cdots\)，\(10\)）. 试验结果如下：

\begin{center}
\begin{tabular}{c|cccccccccc}
\(i\)&1&2&3&4&5&6&7&8&9&10\\\hline
\(x_i\)&545&533&551&522&575&544&541&568&596&548\\
\(y_i\)&536&527&543&530&560&533&522&550&576&536
\end{tabular}
\end{center}

记 \(z_i=x_i-y_i\)（\(i=1\)，\(2\)，\(\cdots\)，\(10\)），记 \(z_1\)，\(z_2\)，\(\cdots\)，\(z_{10}\) 的样本平均数为 \(\bar z\)，样本方差为 \(s^2\).
\begin{enumerate}
\item 求 \(\bar z\)，\(s^2\)；
\item 判断甲工艺处理后的橡胶产品的伸缩率较乙工艺处理后的橡胶产品的伸缩率是否有显著提高（如果 \(\bar z\ge2\sqrt{\frac{s^2}{10}}\)，则认为有显著提高，否则不认为有显著提高）.
\end{enumerate}
\end{problem}

\begin{solution}（1）差值为 \(9\)，\(6\)，\(8\)，\(-8\)，\(15\)，\(11\)，\(19\)，\(18\)，\(20\)，\(12\)，故 \(\bar z=11\)，\(s^2=\frac{1}{10}\sum_{i=1}^{10}(z_i-11)^2=61\).

（2）\(2\sqrt{s^2/10}=2\sqrt{6.1}<11=\bar z\)，所以认为甲工艺处理后的橡胶产品的伸缩率较乙工艺处理后的橡胶产品的伸缩率有显著提高.
\end{solution}

\begin{problem}
记 \(S_n\) 为等差数列 \(\{a_n\}\) 的前 \(n\) 项和，已知 \(a_2=11\)，\(S_{10}=40\).
\begin{enumerate}
\item 求 \(\{a_n\}\) 的通项公式；
\item 求数列 \(\{|a_n|\}\) 的前 \(n\) 项和 \(T_n\).
\end{enumerate}
\end{problem}

\begin{solution}（1）设首项为 \(a_1\)，公差为 \(d\)，则 \(a_1+d=11\)，\(10a_1+45d=40\)，解得 \(a_1=13\)，\(d=-2\)，所以 \(a_n=15-2n\).

（2）当 \(n\le7\) 时，\(T_n=S_n=14n-n^2\)；当 \(n\ge8\) 时，\(T_n=2S_7-S_n=n^2-14n+98\).
\end{solution}

\begin{problem}
如图，在三棱锥 \(P-ABC\) 中，\(AB\perp BC\)，\(AB=2\)，\(BC=2\sqrt2\)，\(PB=PC=\sqrt6\)，\(BP\)，\(AP\)，\(BC\) 的中点分别为 \(D\)，\(E\)，\(O\)，点 \(F\in AC\)，\(BF\perp AO\).
\begin{enumerate}
\item 证明 \(EF\parallel\) 平面 \(ADO\)；
\item 若 \(\angle POF=120^\circ\)，求三棱锥 \(P-ABC\) 的体积.
\end{enumerate}
\bitmapfigure[max width=6.0cm]{img/2023/national_paper_B_liberal/q19_fig1.png}
\end{problem}

\begin{solution}（1）设 \(AF=tAC\)，由 \(BF\perp AO\) 得 \(4(t-1)+4t=0\)，故 \(t=\frac12\)，即 \(F\) 为 \(AC\) 的中点. 于是 \(DE\parallel AB\)，\(OF\parallel AB\)，且 \(DE=OF\)，四边形 \(ODEF\) 为平行四边形，故 \(EF\parallel DO\). 因 \(DO\subset\) 平面 \(ADO\) 且 \(EF\) 不在该平面内，所以 \(EF\parallel\) 平面 \(ADO\).

（2）因 \(PB=PC\)，且 \(O\) 为 \(BC\) 中点，故 \(PO\perp BC\)，从而 \(PO=\sqrt{PB^2-OB^2}=2\). 又 \(OF\parallel AB\perp BC\)，故 \(BC\perp\) 平面 \(POF\). 过 \(P\) 作 \(PM\perp FO\)，则 \(PM\perp\) 平面 \(ABC\)，为三棱锥的高. 由 \(\angle POF=120^\circ\)，得 \(PM=PO\sin60^\circ=\sqrt3\). 故 \(V=\frac13\cdot\frac12\cdot2\cdot2\sqrt2\cdot\sqrt3=\frac{2\sqrt6}{3}\).
\end{solution}

\begin{problem}
已知函数 \(f(x)=\left(\frac1x+a\right)\ln(1+x)\).
\begin{enumerate}
\item 当 \(a=-1\) 时，求曲线 \(y=f(x)\) 在点 \((1,f(1))\) 处的切线方程；
\item 若 \(f(x)\) 在 \((0,+\infty)\) 上单调递增，求 \(a\) 的取值范围.
\end{enumerate}
\end{problem}

\begin{solution}（1）当 \(a=-1\) 时，\(f(1)=0\)，且 \(f'(1)=-\ln2\)，故切线方程为 \(y=-\ln2(x-1)\).

（2）对 \(x>0\)，\(f'(x)=-\dfrac{\ln(1+x)}{x^2}+\dfrac{1/x+a}{1+x}\). 因此 \(f'(x)\ge0\) 等价于 \(a\ge h(x)=\dfrac{(1+x)\ln(1+x)-x}{x^2}\). 由求导得 \(h'(x)<0\)，且 \(\lim\limits_{x\to0^+}h(x)=\frac12\)，\(\lim\limits_{x\to+\infty}h(x)=0\). 故 \(h(x)<\frac12\)，所求范围为 \(a\ge\frac12\).
\end{solution}

\begin{problem}
已知椭圆 \(C:\dfrac{y^2}{a^2}+\dfrac{x^2}{b^2}=1\)（\(a>b>0\)）的离心率是 \(\dfrac{\sqrt5}{3}\)，点 \(A(-2,0)\) 在 \(C\) 上.
\begin{enumerate}
\item 求 \(C\) 的方程；
\item 过点 \((-2,3)\) 的直线交 \(C\) 于 \(P\)，\(Q\) 两点，直线 \(AP\)，\(AQ\) 与 \(y\) 轴的交点分别为 \(M\)，\(N\)，证明：线段 \(MN\) 的中点为定点.
\end{enumerate}
\end{problem}

\begin{solution}（1）由点 \(A\) 在椭圆上得 \(b=2\). 又 \(c/a=\sqrt5/3\)，\(a^2=b^2+c^2\)，解得 \(a=3\)，\(c=\sqrt5\)，故 \(C:\dfrac{y^2}{9}+\dfrac{x^2}{4}=1\).

（2）设直线 \(PQ:y=k(x+2)+3\)，交点横坐标为 \(x_1\)，\(x_2\). 代入椭圆方程得 \((4k^2+9)x^2+8k(2k+3)x+16(k^2+3k)=0\)，从而 \(x_1+x_2=-\dfrac{8k(2k+3)}{4k^2+9}\)，\(x_1x_2=\dfrac{16(k^2+3k)}{4k^2+9}\). 若 \(P(x_1,y_1)\)，\(Q(x_2,y_2)\)，则 \(M\)，\(N\) 的纵坐标分别为 \(\dfrac{2y_1}{x_1+2}\)，\(\dfrac{2y_2}{x_2+2}\). 代入韦达定理化简得其中点纵坐标为 \(3\)，横坐标为 \(0\)，故中点恒为 \((0,3)\).
\end{solution}

\begin{problem}
在直角坐标系 \(xOy\) 中，以坐标原点 \(O\) 为极点，\(x\) 轴正半轴为极轴建立极坐标系，曲线 \(C_1\) 的极坐标方程为 \(\rho=2\sin\theta\left(\frac\pi4\le\theta\le\frac\pi2\right)\)，曲线 \(C_2:\begin{cases}
x=2\cos\alpha,\\
y=2\sin\alpha
\end{cases}\)（\(\alpha\) 为参数，\(\frac\pi2<\alpha<\pi\)）.
\begin{enumerate}
\item 写出 \(C_1\) 的直角坐标方程；
\item 若直线 \(y=x+m\) 与 \(C_1\)，\(C_2\) 均无公共点，求 \(m\) 的取值范围.
\end{enumerate}
\end{problem}

\begin{solution}（1）由极坐标化直角坐标得 \(x^2+y^2=2y\)，即 \(x^2+(y-1)^2=1\)，其中 \(x\in[0,1]\)，\(y\in[1,2]\).

（2）对 \(C_1\)，令 \(\varphi=2\theta\)，则 \(m=y-x=1-\cos\varphi-\sin\varphi\)，其中 \(\varphi\in[\frac\pi2,\pi]\)，故直线与 \(C_1\) 有公共点当且仅当 \(m\in[0,2]\). 对 \(C_2\)，有 \(m=2(\sin\alpha-\cos\alpha)\)，其中 \(\alpha\in(\frac\pi2,\pi)\)，故直线与 \(C_2\) 有公共点当且仅当 \(m\in(2,2\sqrt2]\). 因此两曲线均无公共点时，\(m\in(-\infty,0)\cup(2\sqrt2,+\infty)\).
\end{solution}

\begin{problem}
已知 \(f(x)=2|x|+|x-2|\).
\begin{enumerate}
\item 求不等式 \(f(x)\le6-x\) 的解集；
\item 在直角坐标系 \(xOy\) 中，求由不等式 \(f(x)\le y\) 与 \(x+y-6\le0\) 所确定的平面区域面积.
\end{enumerate}
\end{problem}

\begin{solution}（1）\(f(x)=\begin{cases}
3x-2,&x>2,\\
x+2,&0\le x\le2,\\
-3x+2,&x<0.
\end{cases}\) 分段解不等式得 \(-2\le x\le2\)，故解集为 \([-2,2]\).

（2）平面区域为由折线 \(y=-3x+2\)、\(y=x+2\) 和直线 \(y=6-x\) 围成的四边形. 其左右边界交点横坐标为 \(-2,2\)，面积为 \(\int_{-2}^{0}(4+2x)\,\mathrm dx+\int_0^2(4-2x)\,\mathrm dx=8\).
\end{solution}
